\documentclass[11pt]{article}
\usepackage[margin=1in]{geometry}
\usepackage{amsmath,amssymb,amsthm}
\usepackage[hidelinks]{hyperref}
\newtheorem{theorem}{Theorem}
\newtheorem{corollary}[theorem]{Corollary}
\newcommand{\mr}{\operatorname{mr}}
\newcommand{\rank}{\operatorname{rank}}
\newcommand{\F}{\mathbb F}
\title{Meridional rank from arbitrary permutation branched covers}
\author{Internal research working proof}
\date{17 November 2026}
\begin{document}
\maketitle

A permutation quotient contains information beyond the number of sheets
its meridians move. Keeping the fundamental group of its branched cover
strengthens both the permutation-index bound and the retained restriction
to simple covers. The seven-sheeted cover of the already settled knot
$T(4,7)$ below gives bound four, whereas its permutation index alone and
every simple-cover rank certificate give at most three.
The argument extends to every meridian cycle type. These are written
working proofs and exact internal computations, without a global novelty
claim, formal verification, or external mathematical review. The twenty-four finite equality certificates below require a separate
priority audit before claiming previously unsettled knots. No
counterexample to the Meridional Rank Conjecture is asserted.

\section{The bound for any meridian permutation}

Let $K$ be a classical knot and let
$\rho:\pi_1(S^3\setminus K)\to S_n$, $n\ge2$, have transitive image.
Its connected branched completion is denoted $Y_\rho$.
Suppose a meridian has $c$ cycles on the sheets, including fixed points,
and write $\delta=n-c$. A meridian normally generates the knot group,
so $\delta>0$. Inverting a meridian does not change $c$.

\begin{theorem}\label{thm:general}
For any such representation and any field $k$,
\[
 \mr(K)\ \ge\
 \left\lceil\frac{n-1+\rank\pi_1(Y_\rho)}{\delta}\right\rceil
 \ \ge\
 \left\lceil\frac{n-1+b_1(Y_\rho;k)}{\delta}\right\rceil.
\]
The cover need not be regular, simple, or have involutory meridians.
\end{theorem}
\begin{proof}
Choose any generating family of $d$ meridians and a presentation of the
knot group on these generators. Its one-vertex presentation graph lifts
to a connected graph with $n$ vertices and $dn$ directed edges.
For the $i$th generator, the edges labelled by it are precisely the
$c$ disjoint cycles of its permutation. A cycle of length $a$ represents
the lift of the $a$th power of that generator based at a sheet in the
cycle. This is a meridian of the corresponding lifted boundary torus
which is killed by branched filling. This description remains valid for
any conjugate or inverse of the peripheral meridian.

Add these filling relations for all generators. They are valid in the
branched-cover group; repetitions among them do no harm. Each filling
relation uses only the edges of its one generator cycle. Eliminate one
edge from each cycle using its filling relation. The eliminated edges
are all distinct. Removing an edge of a cycle preserves connectivity,
and no other cycle loses an edge, because these cycles partition the
labelled edges. The graph left after all $dc$ eliminations is therefore
connected with $n$ vertices and $d(n-c)$ edges.

Killing a spanning tree gives $d\delta-n+1$ generators. Lifted knot
relators and any remaining filling relations can only reduce the rank.
Thus $\rank\pi_1(Y_\rho)\le d\delta-n+1$. Minimize over meridional
generating families and use $b_1(Y_\rho;k)\le\rank\pi_1(Y_\rho)$.
\end{proof}

Taking the cover rank to be at least zero recovers the elementary
permutation-index bound $\mr\ge\lceil(n-1)/\delta\rceil$.
When $\delta=1$, the meridian is a single transposition and the theorem
is exactly the retained simple-cover bound.
When $\delta>1$, a small positive cover rank can change the rounded
integer bound even when no simple-cover test detects the knot sharply.
For example, a seven-sheeted representation with double-transposition
meridians gives $\mr\ge4$ as soon as its branched-cover group is
nontrivial. Its permutation-index bound alone is only three.

Define the maximum of the first bound over all finite transitive
permutation representations to be $D(K)$. The maximum exists: its
values are positive integers bounded by $\mr(K)$, and double covers
supply a nonempty set. Then the retained simple-cover invariant obeys
\[
                         C(K)\le D(K)\le\mr(K).
\]
This definition does not give an effective finite bound on the degrees
needed to compute $D(K)$, nor assert equality with $\mr$.

\section{An exact strict comparison on T(4,7)}

Use the torus presentation and meridian convention
\[
 \Gamma=\langle x,y\mid x^4=y^7\rangle,\qquad m=x^{-1}y^2,
\]
already sourced to Cornwell, Section 5.1, in the retained simple-cover
note. On the seven nonzero vectors of $\F_2^3$ take the matrices
\[
 X=\begin{pmatrix}1&0&0\\1&1&0\\0&1&1\end{pmatrix},\qquad
 Y=\begin{pmatrix}0&1&0\\0&0&1\\1&1&0\end{pmatrix}.
\]
They satisfy $X^4=Y^7=(X^{-1}Y^2)^2=I$. Their inverse actions give
right permutation actions, with sheets numbered by the nonzero binary
vectors in increasing binary integer order:
\[
 x=(0\ 6\ 4\ 2)(1\ 5),\qquad
 y=(0\ 2\ 6\ 5\ 4\ 1\ 3),\qquad
 m=(0\ 5)(2\ 3).
\]
The $y$-cycle alone proves transitivity. The meridian has $c=5$ cycles,
so $n=7$ and $\delta=2$.

The lifted graph has fourteen edges. Delete tree edges
$0,1,2,4,8,10$, where edge $7g+v$ starts at sheet $v$ for generator
$g=0,1$. The seven lifts of $x^4y^{-7}$ and one filling word for each
of the five cycles of $m$ give a $12$ by $8$ abelian presentation.
Its Smith factors are
\[
                              1,1,1,1,1,1,1,2.
\]
For a compact independent check, the submatrix in rows
$0,1,3,7,8,9,10,11$ has determinant two, and the binary nullspace of
the full matrix consists of zero and
\[
                             (1,1,0,1,1,0,1,1).
\]
The determinant shows that the presented abelian group is a quotient
of $\mathbb Z/2$; the nonzero binary character proves that quotient
is nontrivial. Thus $H_1(Y_\rho;\mathbb Z)=\mathbb Z/2$.
Theorem~\ref{thm:general} gives
\[
 \mr(T(4,7))\ge\left\lceil\frac{6+1}{2}\right\rceil=4.
\]
The four-strand torus diagram gives the matching bridge upper bound.
This recovers a known torus-knot equality with a stronger cover method.

The retained simple-cover proof establishes $C(T(4,7))=3$: simple
cover degrees are at most five, no onto transposition quotient exists
in degrees three through five, and the double-cover group has rank
two. Consequently
\[
                    C(T(4,7))=3<4=D(T(4,7))=\mr(T(4,7)).
\]
This is a strict comparison of certificate methods, rather than a new
MRC case.

\section{A uniform family with the same strict comparison}

The same strict gap holds for
\[
                             K_j=T(4,420j+7),\qquad j\ge0.
\]
Put $q=420j+7$ and $a=(q+1)/4=105j+2$.
The torus meridian $m=x^{-1}y^a$ has the same seven-sheet permutation,
since $q\equiv0$ and $a\equiv2\pmod7$.
More generally set $q=28t+7$, $a=7t+2$. In the lifted chain matrix
over $\F_2$, replacing $t$ by $t+2$ adds eight complete $y$-cycles to
the torus relation and two complete $y$-cycles to the meridian word.
Each edge then occurs an even additional number of times, with the
same intermediate endpoint. The filling words also have the same
binary edge sums. Hence the matrix is periodic in $t$ modulo two.
The two exact controls $t=0,1$ both have binary nullity one, proving
$b_1(Y_\rho;\F_2)=1$ for every $t\ge0$ and thus $D(K_j)=4$.

For completeness, $C(K_j)=3$ follows without any simple-cover search.
A simple cover has degree at most five, by the four-meridian upper
bound. For $3\le r\le5$, an onto $S_r$ image kills the central element
$x^4=y^q$. Since $\gcd(q,60)=1$, the equation $y^q=1$ in $S_r$
forces $y=1$, so the image is cyclic, a contradiction.
Only the double cover remains. After imposing $m^2=1$, write
$x=y^a m$ and $z=mym$. Its index-two kernel has the presentation
\[
 \langle y,z\mid y^az^ay^az^a=y^q,\quad
                   z^ay^az^ay^a=z^q\rangle.
\]
It has rank at most two. The matrices $Y$ and
$Z=(X^{-1}Y^2)Y(X^{-1}Y^2)$ satisfy these two relators and do not
commute, so its rank is exactly two. Therefore
\[
                       C(K_j)=3<4=D(K_j)=\mr(K_j)=b(K_j).
\]
All these torus-knot equalities are already known. This family proves
that admitting more general meridian cycle types removes a uniform
limitation of the simple-cover invariant.

\section{Twenty-four sharp four-meridian certificates}

The retained finite witness input contains twenty-nine thirteen-crossing
knot diagrams with natural and dual seven-sheeted $GL_3(\F_2)$ actions.
For the following twenty-four names, the corresponding branched-cover
groups are cyclic of order three for both actions:
\[
\begin{array}{rrrr}
13a_{373}&13a_{375}&13a_{391}&13a_{392}\\
13a_{3264}&13a_{3266}&13a_{3275}&13a_{3276}\\
13n_{1005}&13n_{1006}&13n_{1007}&13n_{1011}\\
13n_{1012}&13n_{1013}&13n_{1043}&13n_{1044}\\
13n_{1045}&13n_{1046}&13n_{1047}&13n_{1048}\\
13n_{3842}&13n_{3844}&13n_{3855}&13n_{3856}.
\end{array}
\]

\begin{theorem}\label{thm:twentyfour}
For every one of these twenty-four knots, $\mr(K)=b(K)=4$.
\end{theorem}
\begin{proof}
The explicit matrices and literal planar diagrams give transitive
seven-sheeted actions, each meridian a product of two disjoint
transpositions. Thus $n=7$ and $\delta=2$.
The exact cover presentations described below establish
$\rank\pi_1(Y_\rho)=1$. Theorem~\ref{thm:general} gives
$\mr(K)\ge\lceil(6+1)/2\rceil=4$.
The four-seed propagation certificates on the same literal diagrams
provide the matching bridge upper bounds: the Wirtinger number is at
most four, and it equals bridge number by Blair--Kjuchukova--Velazquez--
Villanueva \cite[Theorem 1.3]{wirt}. The separate diagram
checker verifies their full propagation and planarity; it does not
assume bridge number four from an open-list interval.
\end{proof}

Here is the group-presentation certificate. For each action, every
fixed-sheet lifted meridian edge dies under filling, and the two edges
of each moved pair become inverses. There are twenty-six undirected
edges on seven vertices. Delete a spanning tree of six edges, leaving
twenty generators. Lift every diagram crossing relation from all seven
sheets. Since all filled meridian paths are involutory, a crossing
relation becomes the appropriate four-step path $o\,a\,o\,b$,
independently of its crossing sign. These are exact presentations of
the branched-cover groups, rather than presentations of chosen finite
quotients of those groups.

For each of the forty-eight positive actions, nineteen explicit
one-occurrence Tietze eliminations leave the presentation
$\langle g\mid g^3\rangle$. At every step, the defining relation
contains the eliminated generator exactly once, so solving for it
and substituting everywhere preserves the group. A small verifier
reconstructs the folded presentation and replays every substitution.
The separate full-directed-edge homology checker gives
$H_1(Y_\rho;\mathbb Z)=\mathbb Z/3$ independently of these group
simplifications. Neither calculation assumes that cyclic homology
implies a cyclic group.

For each natural and dual action on the five other retained knots,
\[
             13n_{4633},\quad13n_{4634},\quad13n_{4637},\quad
             13n_{4638},\quad13n_{4643},
\]
twenty one-occurrence Tietze eliminations remove all generators and
relations. Their ten cover groups are therefore trivial. The best
bound from each of these fixed actions is three, even when full cover
group rank replaces first homology. This excludes a nonlinear
finite-quotient improvement for these particular covers; it does not
exclude other permutation actions or establish any knot's meridional
rank to be three.

The literal diagrams, meridian matrices and seed traces suffice to
replay these certificates. Their Hoste--Thistlethwaite identities are
checked by transfer to the independently retained named census data;
a diagram label in a search result is not taken as identity evidence.
All twenty-nine names occur in the author repository still-open list
\cite{bhattacharjee}, checked on 17 November 2026 at Git blob
\texttt{01c7c332a0bf1565d7558aec4c769ff36cbd31a9}. This is evidence
that the names were listed there, not an assertion that no earlier
theorem covers them. No completeness of the retired broad
$GL_3$ search is needed for the twenty-four positive statements.
The assertion is a finite equality theorem with exact current
certificates. Its relationship to previously established knot families
and worldwide priority requires a separate source audit; the count of
names alone does not establish twenty-four new MRC cases.

\section{Verification scope and next use}

The compact checker reconstructs the right actions directly from the
binary matrices, verifies transitivity and every torus/peripheral word,
constructs the full lifted edge matrix, and proves the displayed integral
homology by an exact determinant and exhaustive binary character check.
The initial Smith calculation is not needed as a trusted checker input.
The second parity control and the two-generator double-cover matrix
relations are checked separately. Deliberate corruptions of a relation and a sheet permutation, and a
nontransitive action, are rejected. Omitting fixed-sheet fillings changes
the checked homology, detecting that essential boundary error.
These finite controls support the stated examples; the uniform theorem
and parity argument are written proofs.

The retained $GL_3(\F_2)$ knot quotients whose
meridians are transvections supply the preceding finite application. Their natural seven-point actions have
$\delta=2$. Any nontrivial branched-cover group gives a sharp lower
bound four when a four-meridian upper certificate is available.
Failure of the earlier rank-one cocycle test on the three-dimensional
module does not exclude this different permutation-module or nonlinear
cover-group test. Conversely, trivial first homology does not prove
that a cover group is trivial.
\section{Classical context}

Homology calculations for arbitrary finite branched covers have a
classical Fox-calculus framework. Hempel \cite[Section 3, Theorems
3.1--3.2]{hempel} computes first homology using relative Fox matrices,
and describes the meridian-cycle filling data on printed page 121.
The lifted-edge matrices here implement that classical construction.
The exact meridional-generator inequality and strict comparison above
are proved directly, without attributing an unlocated stronger
statement to that source. Their worldwide priority remains unclaimed.

\begin{thebibliography}{9}
\bibitem{wirt} R. Blair, A. Kjuchukova, R. Velazquez and P. Villanueva,
Wirtinger systems of generators of knot groups,
Communications in Analysis and Geometry 28(2) (2020), 243--262.
\url{https://doi.org/10.4310/CAG.2020.v28.n2.a2}.
\bibitem{bhattacharjee} D. Bhattacharjee, author repository
\emph{conjecture-to-choose}, \texttt{data/still\_open.json}, checked
17 November 2026, exact Git blob given in the text.
\url{https://github.com/anaxia-research/conjecture-to-choose/blob/main/data/still_open.json}.
\bibitem{hempel} J. Hempel, Homology of branched coverings of
3-manifolds, Canadian Journal of Mathematics 44(1) (1992), 119--134.
\url{https://doi.org/10.4153/CJM-1992-007-4}.
\end{thebibliography}
\end{document}
